% a02_Kpos part deg1: the Ward rows of the degree count on the f = 1 flow (input from a02_Kpos.tex after thm:a02:T)
\subsection{The Ward rows of the degree count}\label{ssec:a02:ward}

The source-free, block-free classes of King's expansion with $D\le0$ are, by order and external profile (\cite[(3.30)--(3.31)]{K1}: $D=\frac o2-3+\frac12(E_S+E_A)+\frac32E_{S_D}$ for a member of order $o$ with $E_S$ undifferentiated scalar, $E_{S_D}$ differentiated scalar and $E_A$ photon legs), the vacuum classes, the photon one- and three-point classes, the scalar two-point classes with and without one photon leg, the photon two-point classes of orders $e^2$, $e^4$, $e^2\lambda$, and the tree \cite[(2.4)]{Ba3}. The resummed coefficient $c_S$ of such a class vanishes, or is coarse-dominated, by one of six mechanisms: the counterterm matched per scale assignment \cite[(3.33)--(3.35)]{K1}; charge parity; a derivative at the end of a line; the locality of the tadpoles; a reflection; and the zero-momentum gauge covariance of the full fixed-order two-photon function and of the photon insertions into the exchange graph $\Sigma_1$. The rows of the last kind, the photon two-point classes (all $24$ configurations, with either photon leg an $A'$-line or an $\tilde A$ field) and the scalar two-point classes with one photon inserted into $\Sigma_1$ at a vertex or on its scalar line, or into the $\lambda$-tadpole loop, are the \emph{Ward rows}. They differ from the others in two respects. Their identity holds for the \emph{whole inside family} with all internal lines resummed and replaced by one kernel, not member by member or at a fixed index assignment; so at a fixed localisation tuple the localisation weights must be common to all members of the family, which under King's rule of one partition of unity per internal vertex \cite[(3.36)]{K1} they are not (at order $e^2$ the bubble has two internal vertices and the $J_2$-loop one, so at a fixed tuple the bubble has no partner, and its weighted zeroth moment grows like $\xi^{-1}$). And in the $R>p$ branch \cite[(3.38)]{K1} King cuts the lines longer than $3\sqrt d$ by vertex positions \cite[p.~666]{K1}, which the identities sum over. Both are met by a reading of King's localisation and cut on data fixed by the tuple alone; this subsection gives it, on the $f=1$ flow, and proves the coarse domination $|c_a|\le C\rho_{j_0}^{D(S)+|a|}$ of \cref{thm:a02:T}(ii) on the Ward rows. Throughout, $S$ is a prefix component of an ordering with $D(S)\le0$, no source, block or $h$-vertex, legs $\Phi_e$ (external fields or lines of index $\ge j_0$ evaluated as fields) of indices $j_e$, $j_0\defeq\min_ej_e$; its \emph{inside family} $F_S$ is the set of all connected block-free subgraphs of the same order and profile occurring in place of $S$ with the same legs at the same outside points, the $\dm$-vertex and $\dm$-inserted graphs included \cite[p.~431]{Ba3}; the \emph{admitted set} $A_S$ is all internal indices $\le j_0$; and for a reference vertex $v_0$ (the vertex of the coarsest leg) and a multi-index $a$ on the other legs the moments are $c_a(x)\defeq\sum_{\text{internal positions}}5^{-3j(V-1)}K_{F_S}(x;x_2,\dots)\prod_{e\ge2}(x_e-x)^{a_e}/a!$, $x$ the position of $v_0$, internal localisations summed.

\begin{lemma}[resummation]\label{lem:a02:RS}
For fixed outside indices, the sum over all orderings of the lines of $S$ and all index assignments in $A_S$ equals the expression with every line of $S$ and of every member of $F_S$ (the internal lines of $\dm$ units included) replaced by the resummed line $G_{j_0}=\sum_{i\le j_0}G_{(i)}$ (resp.\ $G^{\mathrm{ph}}_{j_0}$). Rescaled by $\rho_{j_0}$ to the lattice $\xi\Z^3$, $\xi\defeq5^{-j_0}$, let $C^\xi\defeq(-\Delta+1)^{-1}$ (scalar $\otimes1$; $\delta_{\mu\nu}$ for the photon) and let $F^C_S$ be the family with every line, counterterm-internal lines included, replaced by $C^\xi$ on $\xi\Z^3$. Then for $|a|\le n\defeq-D(S)$
\[
c_a(F_S)=\rho_{j_0}^{D(S)+|a|}\Big[c_a(F^C_S)+\sum_{\emptyset\ne T}c_a(F_{S,T})\Big],\qquad|c_a(F_{S,T})|\le C\ \text{uniformly in }\xi,x,j_0,
\]
where $F_{S,T}$ has the lines of $T$ replaced by $\Delta\defeq G_{j_0}-C^\xi$ (rescaled) or by the periodisation tail on a torus, and the others by $C^\xi$.
\end{lemma}
\begin{proof}
The first statement is \cite[(3.57)--(3.58)]{K1} and multilinearity: King's tie rule is replaced by a simultaneous $S$-first tie-break, which changes nothing in the total. For the second, the resolvent identity $G=C^\xi+G(1-\tilde m^2-\tilde aP)C^\xi$ \cite[p.~437]{Ba3} ($P=Q^*Q$, $\tilde m^2=m^2l_{j_0}^2$, $\tilde a=a_{j_0}$) gives $\Delta=G(1-\tilde m^2-\tilde aP)C^\xi$; with $|G(x,y)|\le C|x-y|_\xi^{-1}e^{-\delta|x-y|}$, $|DG|\le C|x-y|_\xi^{-2}e^{-\delta|x-y|}$, $|D_xD_yG|\le C|x-y|_\xi^{-3}e^{-\delta|x-y|}$ (\cref{prop:a02:FL}; \cite[Thm.~3.3]{K1}), the same for $C^\xi$, and $P$ bounded with unit-range kernel, convolution in $d=3$ gives $|\Delta|\le C$, $|D\Delta|\le C(1+\ln_+(1/|x-y|_\xi))$, $|D_xD_y\Delta|\le C|x-y|_\xi^{-1}$, all times $e^{-\delta'|x-y|}$: a $\Delta$-line is a unit-scale kernel, undifferentiated bounded and Lipschitz below scale $1$, differentiated two orders better than the $G$-line (degree $+2$ up to a logarithm absorbed by an $\epsilon$-loss). The replacement is multilinear, the $\dm$ partners' internal lines being expanded by the same rule, so every all-$C$ subgraph of $F_{S,T}$ is renormalised inside $F_{S,T}$ by its own partners. On a torus $C^\xi_T=\sum_\gamma C^\xi(\cdot+\gamma)$ and the $\gamma\ne0$ tail is smooth and $O(e^{-\delta|\gamma|})$, one more $\Delta$-type line. Decompose the $C^\xi$ lines into scales ($C^\xi$ is the $A=0$, mass-$1$ instance of \cite[(2.13)--(2.17)]{K1} on $\xi\Z^3$ with the bounds \cite[(3.63)--(3.65)]{K1}) and order all lines: a prefix component of an output term is either all-$C$, hence a proper subgraph of a member of $F_S$ (lower order, or the same order and more legs), whose class is rewritten by the same statement for the $C^\xi$ line system (induction on order, then on the number of legs, which is well founded since $E\le4$ for every divergent class); or it contains a $\Delta$-line: undifferentiated, it sits at the top scale and enters no fine-scale sum; differentiated, its degree is $\ge D_{\mathrm{plain}}+2-\epsilon\ge1-\epsilon>0$, because every non-vanishing non-vacuum source-free connected subgraph has $D\ge-1$ (the $D=-2$ classes are the current tadpole, which vanishes by charge parity, and vacuum graphs, which are not prefix components of a non-vacuum $S$). So every ordered subgraph of every output term has positive degree, the multiscale bound \cref{thm:a02:S3}(a) applies, and the final scale sum is coarse-dominated.
\end{proof}

\begin{lemma}[zero-momentum gauge covariance]\label{lem:a02:GC}
Let the lines be gauge covariant in a background $B$, $C_B(x,y)\mapsto g(x)C_B(x,y)g(y)^T$ under $B\mapsto B+\nabla\chi$, $g=e^{eq\chi}$, with King's vertices \cite[(3.18)--(3.23)]{K1} at background $B$, and on $\xi\Z^3$ let $B$ be constant, $\chi_B(x)=B\cdot x$. For a labelled contraction $G$ with two external scalar fields at $x,x_2$, $K^G_B(x,x_2)=g(x)K^G_0(x,x_2)g(x_2)^T$, and for $G$ with photon legs only $F^G_B=F^G_0$. Differentiating at $B=0$ in direction $\mu$: \textup{(a)} with $W(A)\defeq\ln Z(A)-\ln Z(0)$, $Z(A)=\int d\mu(A_{\mathrm{int}})\int d\mu_{C_0}(\phi)e^{V(\phi,A+A_{\mathrm{int}})}$, $W(A+B)=W(A)$ for constant $B$ on $\xi\Z^3$, hence for the full fixed-order connected two-photon function $\Pi=\partial^2W/\partial A\partial A|_0$, all its labelled contractions and $\dm$-inserted ones included, $\sum_{b_2\in\mathrm{bonds}_\nu}\Pi_{\mu\nu}(b_1,b_2)=\partial_{A_{b_1}}\partial_{B_\nu}W(A+B)|_0=0$; \textup{(b)} for every two-scalar-leg parent $G$ with scalar charge structure, $\sum_{b\in\mathrm{bonds}_\mu}K^{(G+\mu)}(x,x_2;b)=eq(x-x_2)_\mu K^G_0(x,x_2)$, so the zeroth moment of the insertion family of $G$ equals $-eq$ times the first moment of $G$. The insertions are: on a scalar line $-C(\partial_B(-\Delta_B))C$, the vertex $J_1$; at a vertex $J_m(B)$ through $U(B_b)$ a $J_{m+1}$-type vertex and through $D_B$ a $Q_{m+1}$-type vertex; none on photon lines, $\lambda$ or $M$.
\end{lemma}
\begin{proof}
(a) $Z(A+\nabla\chi)=Z(A)$ by the change of variables $\phi\mapsto e^{eq\chi}\phi$, a constant $B$ is $\nabla(B\cdot x)$ on $\xi\Z^3$, and the order-by-order expansion of $W$ in $e$ commutes with the shift because $W$ depends on $eA$ only. (b) Per contraction the value is a contraction of covariant tensors with covariant propagators, the internal phases cancel, $\partial_{B_\mu}g=eqx_\mu g$, and $K^G_0$ commutes with $q$; the insertion list is Leibniz on $C_B$ and on $J_m(B)$ \cite[(3.22)]{K1} ($Q_m$ carries no $U$ and is $B$-independent).
\end{proof}
The identity is used only for $F^C_S$ of \cref{lem:a02:RS}, whose lines are full resummed lines replaced by $C^\xi$ (the total-field form); no identity is applied to a fixed ordering or a fixed index assignment. On a torus (a), (b) hold up to the holonomy response $O(e^{-\delta|T|_\xi})$, one more $\Delta$-type term.

\begin{lemma}[reflection]\label{lem:a02:REF}
Let $g_s$, $c$ be the even scalar and photon kernels $C^\xi$ on $\xi\Z^3$. \textup{(a)} The exchange kernel $\Sigma_1$ with both $S_D$ legs on the scalar line and $\mu=\nu$ is $k(y)=-e^2\sum_\mu(\Delta_\mu g_s)(y)c(y)$, $y=x-x_2$, which is even, so its first moment vanishes. \textup{(b)} For the order-$e^2$ two-photon family ($J_1J_1$ bubble, $J_2$ loop, $Q_2$ loop) with zeroth moment $M^{(0)}=0$, every first moment $M^{(1)}_{\mu\nu\sigma}\defeq\sum_{b_2}\Pi_{\mu\nu}(b_1,b_2)(b_2-b_1)_\sigma$ vanishes.
\end{lemma}
\begin{proof}
(a) With $b=(x-\xi e_\mu,x)$, $b'=(x_2-\xi e_\mu,x_2)$, $\partial_b\partial_{b'}g_s(x-x_2)=-\Delta_\mu g_s(x-x_2)$ and the photon line is $c(x-x_2)$. (b) Under the reflection $R_\sigma$ the bond $(z,\mu)$ goes to $(Rz-\xi e_\sigma\delta_{\mu\sigma},\mu)$ and $A_\mu\mapsto s_\mu A_\mu$ ($s_\sigma=-1$, else $+1$); the family is the expansion of a reflection-invariant action, so $\Pi_{\mu\nu}(b_1,b_2)=s_\mu s_\nu\Pi_{\mu\nu}(R_\mu b_1,R_\nu b_2)$, hence $M^{(1)}_{\mu\nu\sigma}=s_\mu s_\nu[-M^{(1)}_{\mu\nu\sigma}-\xi(\delta_{\nu\sigma}-\delta_{\mu\sigma})M^{(0)}_{\mu\nu}]$ under $R_\sigma$ and $=s'_\mu s'_\nu M^{(1)}_{\mu\nu\sigma}$ under $R_\tau$, $\tau\ne\sigma$. With $M^{(0)}=0$: if $s_\mu s_\nu=+1$ under $R_\sigma$ then $M^{(1)}=0$; otherwise exactly one of $\mu,\nu$ is $\sigma$, say $\mu=\sigma\ne\nu$, and $R_\nu$ gives $M^{(1)}=-M^{(1)}$.
\end{proof}

\paragraph{The localisation by centres.} King localises every graph of $P^{(j),1}$ by inserting $1=\sum_z\chi_z$ at each internal vertex \cite[(3.36)]{K1} and summing over the tuple $(\{y_i\},\{z\})$ \cite[(3.37)]{K1}. We attach the partitions of unity to \emph{centres} instead: \textup{(i)} one centre to every $A'$ field of a $J_m$ or $Q_m$ vertex ($m$ centres at $J_m$ \cite[(3.22)]{K1}, $m$ at $Q_m$ \cite[(3.23)]{K1}); \textup{(ii)} one centre to every other vertex ($\lambda$ \cite[(3.18)]{K1}, block vertices \cite[(3.24)--(3.26)]{K1}, $h$-, $g$- and source vertices). For each centre $c$ insert $1=\sum_{z_c}\chi_{z_c}(p_c)$, $p_c$ the position of the vertex carrying $c$; the tuple of a graph is $\tau=(\{y_i\},(z_c)_c)$; $\square$ is the smallest cube of large blocks containing every $y_i$ and every $z_c$ (the centres of the localised counterterm graphs included), and $\square'$, $\square''$, $x_0$, $B=A_j(x_0)$, $\tilde A$ and $R\defeq\dist(\{y_i\},\{z_c\})$ are as in \cite[pp.~660--662]{K1}. King's index rules \cite[p.~667]{K1} are kept: a $\dm$ graph is localised only if all its integers are $<j$, and an $E_1$ piece with $m_1<m$ gets one centre at one vertex connected to $l(m)$. King's printed rule is this rule with one centre per vertex in place of (i).

\begin{lemma}[King's localisation argument with centres]\label{lem:a02:Lsharp}
\textup{(a)} The only properties of \cite[(3.36)]{K1} used in \cite[Sects.~3.4--3.5]{K1} are: $\sum_z\chi_z=1$ per inserted factor; $0\le\chi\le1$ with support of size $O(1)$ and bounded derivatives; every localised vertex lies within $O(1)$ of a centre in $\square$; and the number of centres is bounded by a function of $\nbar$ and the order. The rule by centres has all of them, so \cite[Thm.~3.5, Prop.~3.6, (3.40)--(3.42), (3.43)--(3.57), (3.68), (3.85)]{K1} and the step comparison of \cite[p.~665]{K1} hold verbatim, with constants depending on $\nbar$ and $d$ (the centres of one vertex lie within $\frac43\sqrt d$ of each other, so the tree distance over centres exceeds that over vertices by at most $C(\nbar)$). \textup{(b)} \emph{Centre count.} For $S'\in F_S$ with $e$-order $\alpha$, $\lambda$-order $\beta$ and $E_{\tilde A}$ legs that are $\tilde A$ fields, the number of centres is $N(S')=\alpha-E_{\tilde A}+\beta$; it is the same for every member of $F_S$. \textup{(c)} Consequently every graph of a renormalisation orbit carries the same tuple, under a fixed bijection of the centres of $S$ and $S'$, hence the same $\square$, $\square'$, $x_0$, $B$, $\tilde A$ and the same branch of \cite[(3.38)/(3.39)]{K1}; the same along any sequence of replacements, nested units included. \textup{(d)} In the order-$e^2$ photon two-point family every centre sits at the $S$-end of an $A'$-line leg: no member has an internal photon line or a $\lambda$ vertex, and two $\tilde A$ legs are impossible.
\end{lemma}
\begin{proof}
(a) Each inserted factor sums to $1$, so \cite[(3.37)]{K1} is an identity for any number of centres; the uses on \cite[pp.~659--668]{K1} are: the definition of the tree distance and the decay $\exp[-\delta\dist]$ (co-located centres cost a bounded factor); the dichotomy and $\sum_{y,z}e^{-\frac12\delta R}\le C\eta^{-3}|\CJ|$; the cubes; $B=A_j(x_0)$ on $\square'$; the gauge transformation and $\theta$; $|x_0-z_c|\le Cp$ by the tree decay; \cite[(3.68)]{K1}, the weight of a vertex being a product of $\chi$'s in $[0,1]$ supported in a unit cube; the root sum by $|\square'|$; the replacement of $\chi(x')$ by $\chi(x)$ in the step comparison, since products of $\chi$'s have bounded derivatives; $\delta G$ bounds ``since $x,y\in\square$'', every vertex lying within $\frac23\sqrt d$ of each of its centres; the counterterm localisation with the index rules; and the fold \cite[(3.85)]{K1} of a leg's $\chi$ into the leg. (b) Each $A'$ field at a $J$- or $Q$-vertex carries one factor $e_j$ \cite[(3.22)--(3.23)]{K1}; each $\tilde A$ carries one $e_j$ \cite[(3.47)--(3.53)]{K1} and is never contracted; each $\lambda$ vertex carries one $\lambda_j$. So $\alpha=\sum_v(a_v+t_v)$ with $a_v$ the $A'$ fields and $t_v$ the $\tilde A$ legs at $v$, $E_{\tilde A}=\sum_vt_v$, $\beta=|\Lambda(S')|$, and $N=\sum_va_v+|\Lambda(S')|$. $F_S$ fixes $(\alpha,\beta)$ and the legs; every identity used acts within one $(\alpha,\beta)$, since gauge invariance holds for every value of $e$ and $\lambda$ and $\dm$, $E_1$ are sums of graphs of definite order \cite[(3.33)--(3.35)]{K1}. (c) The reduction is graph-local and the listed objects are functions of the tuple. (d) $\alpha=2$, $\beta=0$, $E_{A'}+E_{\tilde A}=2$, so there are exactly $E_{A'}$ $A'$ fields, the $S$-ends of the $A'$-line legs; a $J$/$Q$ vertex needs an $A'$ field, and a subgraph of $\tilde A$-only vertices is the tree \cite[(2.4)]{Ba3}, not a two-photon graph.
\end{proof}

\begin{lemma}[frozen weights]\label{lem:a02:FW}
Let $S$ be a Ward-row unit at a fixed tuple, $x$ its reference point, and for a member $S'$ let $W_{S'}=\prod_{c\in\mathrm{legC}}\chi_{z_c}(p_c)\prod_{c\in\mathrm{intC}}\chi_{z_c}(p_c)$, where $\mathrm{legC}$ are the centres at the $S$-ends of the $A'$-line legs (common to all members) and $\mathrm{intC}$ the remaining $N-E_{A'}$ centres (common by \cref{lem:a02:Lsharp}(b), identified by the fixed bijection). Put $\tilde\Phi_\ell\defeq\chi_{z_\ell}\Phi_\ell$ for each $A'$-line leg, $\tilde\Phi\defeq\Phi$ for the other legs, and $W^\circ\defeq\prod_{c\in\mathrm{intC}}\chi_{z_c}(x)$, a member-independent number. Then
\[
\sum_{S'\in F_S}\sum_{\mathrm{pos}}K_{S'}W_{S'}\prod_\ell\Phi_\ell=W^\circ\sum_{S'}\sum_{\mathrm{pos}}K_{S'}\prod_\ell\tilde\Phi_\ell+\sum_{S'}\sum_{\mathrm{pos}}K_{S'}\Big[\prod_{c\in\mathrm{intC}}\chi_{z_c}(p_c)-W^\circ\Big]\prod_\ell\tilde\Phi_\ell.
\]
The first term is the unweighted family sum against the legs $\tilde\Phi$ (unit-scale smooth factors times legs, H\"older or $C^{1,\alpha}$ at the leg's own scale, \cite[(3.85), p.~665]{K1}), times a constant. The second term is a sum of remainder terms with a credit $\rho_{j_l}$ on a line $l$ of $S$ and a debit on a unit-scale weight (cost $\rho_j^{-1}=1$), so $D'(S)=D(S)+1$ and $D'(\gamma)\ge D(\gamma)$ for $\gamma\subsetneq S$; it is absent for the order-$e^2$ photon two-point class ($\mathrm{intC}=\emptyset$ by \cref{lem:a02:Lsharp}(d)), the only Ward row with $D=-1$, where the $n=1$ Taylor step acts on $\tilde\Phi$, King's fold \cite[(3.85)]{K1}.
\end{lemma}
\begin{proof}
Add and subtract; for products of $[0,1]$-valued functions $|\prod a_i-\prod b_i|\le\sum|a_i-b_i|$, so $|\prod\chi(p_c)-W^\circ|\le\sup|\nabla\chi|\sum_{c}|p_c-x|\le C\sum_{l\in P}|\mathrm{step}_l|$ along a path $P$ of lines of $S'$ from $x$ to the vertices carrying $\mathrm{intC}$; each step is absorbed by part of its line's decay at the price $\rho_{j_l}$ \cite[(3.13)--(3.14), (3.19)]{Ba3}. Every Ward row other than the order-$e^2$ photon class has $D(S)=0$, hence $D'(S)=1>0$.
\end{proof}
